\section{Introduction}

\subsection{Problem Statement and Motivation}
% [REVISION: Moved Clinical Problem to the top for better narrative flow]
Epilepsy affects roughly 7.6 per 1,000 people, and approximately 30\% of
patients are drug-resistant, motivating a critical demand for timely,
low-burden monitoring and early intervention \parencite{beghiEpidemiologyEpilepsy2019,
kwanEarlyIdentificationRefractory2000, chenNewEraPersonalised2020}.
While EEG remains the reference standard in controlled settings, EEG
hardware is obtrusive for pervasive use. Wearables offer the possibility
of continuous monitoring in ambulatory settings, addressing the need for
safety and better seizure management \parencite{chenNewEraPersonalised2020}.

Artificial intelligence (AI) and, in particular, deep learning have
enabled substantial advances in medical signal analysis, supporting
early detection and continuous monitoring. The focus of this literature
review is on time-series deep learning models (CNN, LSTM, TCN) applied
to wearable seizure detection and prediction using non-EEG wearable
and autonomic biosignals (ECG/HRV, PPG, EDA, ACC).

Healthcare constraints sharpen these modeling objectives. Models must
achieve high sensitivity while strictly controlling false alarms per
hour (FAR/h) to avoid alarm fatigue. They must also balance predictive
performance with the interpretability and safety requirements demanded
in regulated clinical contexts. Furthermore, robustness to domain shift
across patients, devices, sites, and activities is required.
Consequently, evaluations must address data scarcity by prioritizing
prospective, multimodal, and patient-independent study designs.

% [REVISION: Clarified the operationalization]
The generalized topic of time-series deep learning models is
operationalized here via a focused medical use case: comparing
traditional feature-based models against emerging deep temporal
architectures (1D-CNN, LSTM, TCN, attention/transformers) for seizure
detection and preictal prediction. The key question is determining
which modeling and evaluation choices deliver clinically viable, low
FAR/h, high-sensitivity performance under pseudo-prospective and
prospective ambulatory conditions.

\subsection{Research Aim and Questions}
This literature review synthesizes machine and deep learning models for
seizure detection and prediction using non-EEG wearable/autonomic
biosignals to determine optimal strategies for achieving clinically
meaningful performance in ambulatory contexts.

\subsubsection{Primary Research Question}
In non-EEG wearable seizure detection and preictal prediction
(ECG/HRV, PPG, EDA, ACC), which modeling and evaluation choices
reported in 2015--2025 are associated with low alarm-based
false-alarm rates and high sensitivity with acceptable latency in
ambulatory settings?

\subsubsection{Sub-questions}
\begin{enumerate}
  \item What ranges of sensitivity, FAR/h, and latency are reported under
    retrospective, pseudo-prospective, and prospective designs, and how
    do alarm-based results differ from sample-based metrics?
  \item Across studies, how do traditional feature-based models compare
    to 1D-CNN, TCN, and LSTM architectures, and what evidence exists for
    attention or transformer variants on these signals regarding
    computational and latency considerations?
  \item How do temporal choices (window length, stride, context length)
    correlate with reported alarm-based performance in real-time use?
  \item What evidence is available on robustness to domain shift
    (patient, device, site, activity), including cross-dataset or
    external validation, and how do these factors affect alarm-based
    metrics?
  \item What dataset gaps and class imbalance issues are reported, and
    what is the measured impact of augmentation or synthetic data on
    alarm-based outcomes?
\end{enumerate}

\medskip
\noindent\textit{Contribution.}
This review contributes a systematic synthesis of deep learning
architectures specifically for non-EEG seizure forecasting. Unlike
previous reviews, it explicitly analyzes the trade-off between
sensitivity and FAR/h in ambulatory deployments, offering practical
recommendations for model design and evaluation protocols to bridge
the gap between research prototypes and clinical utility.

\subsection{Methodological Approach and Structure}
The review follows structured IS and medical literature guidance
(concept matrix inspired by \cite{websterAnalyzingPastDetermine2002}
and PRISMA principles). The methodological approach uses a multi-stage
search and transparent extraction/synthesis procedures as outlined
below.

\begin{enumerate}
  \item \textbf{Search strategy (Multi-stage, 2015--2025):}
    \begin{itemize}
      \item \textit{Scoping:} Google Scholar for broad coverage and
        identifying candidate papers.
      \item \textit{Formal queries:} IEEE Xplore, PubMed, and Scopus
        using controlled keyword strings and deduplication.
      \item \textit{Citation chasing:} Backward and forward citation
        tracking on key papers to ensure completeness.
    \end{itemize}
  \item \textbf{Inclusion / Exclusion criteria:}
    \begin{itemize}
      \item \textit{Inclusion:} Studies using wearable/non-EEG autonomic
        signals for seizure detection or preictal prediction.
      \item \textit{Exclusion:} EEG-only studies, invasive recordings,
        or papers limited to purely algorithmic simulations without
        human data.
    \end{itemize}
\end{enumerate}

Extracted dimensions populate a concept matrix with the following
fields:
\begin{enumerate}
  \item Setting (clinical vs. ambulatory).
  \item Modality (ECG/HRV, PPG, EDA, ACC).
  \item Task (detection vs. prediction).
  \item Model class and fusion strategy.
  \item Study phase (retrospective, pseudo-prospective, prospective).
  \item Evaluation granularity (sample- vs. alarm-based).
  \item Windowing specifics (window length, stride, overlap).
  \item Metrics (sensitivity, FAR/h, AUC, F1-score, latency).
  \item Personalization strategies.
  \item Interpretability and safety considerations.
  \item Data characteristics and dataset identity.
  \item Missing data handling and synthetic data use.
  \item Inference requirements (latency, computational load,
    hardware).
\end{enumerate}

\noindent\textbf{Prioritization Criteria:}
Studies utilizing patient-independent splits, reporting FAR/h, and
conducting external or cross-dataset validation will be prioritized in
the synthesis. A PRISMA-style flow diagram will summarize and visualize
the screening process.

% [REVISION: Moved Background BEFORE Modeling Landscape]
\section{Theoretical Background and Concepts}

\subsection{Epileptic Seizures}
Epileptic seizures are classified by their onset as focal (starting in
a specific area of the brain) or generalized (affecting both
hemispheres from the onset). Focal seizures are further categorized
into focal aware and focal impaired awareness seizures, depending on
patient awareness, and by onset features (motor vs. non-motor). Focal
seizures may evolve to bilateral tonic–clonic seizures. Generalized
seizures begin in both cerebral hemispheres and include motor (e.g.,
tonic–clonic, myoclonic, atonic) and non-motor (absence) types
\parencite{fisherInstructionManualILAE2017}. Tonic–clonic seizures
commonly cause loss of consciousness and carry increased injury and
SUDEP risk \parencite{hesdorfferCombinedAnalysisRisk2011}.

\subsection{Autonomic Markers and Wearable Sensing}
Sympathetic activation around seizures manifests as ictal tachycardia,
preictal HR rise with reduced HRV (lower RMSSD/HF, higher LF/HF),
increased EDA (tonic SCL, phasic SCRs), and limited ACC changes
useful for artifact removal and context awareness. PPG yields
pulse-rate variability analogous to HRV from ECG
\parencite{mironAutonomicBiosignalsSeizure2025,masonHeartRateVariability2024}.

Wearable heart-rate sensing aligns closely with chest-strap ECG under
motion (e.g., Apple Watch Ultra/Polar H10 concordance $\sim$0.998),
though dominant-wrist motion degrades accuracy. Chest-strap ECG
provides robust RR intervals during exercise
\parencite{wolfHearttoWearAssessingAccuracy2024a}.

% [REVISION: Broke this huge block of text into subsections for readability]
\section{Modeling Landscape and Design Dimensions}

Detection and preictal prediction from ECG/HRV, PPG, EDA, and ACC are
framed by windowing, labeling, and alarm logic, with alarm-based
sensitivity, FAR/h, and latency as the outcomes used in the research
questions.

\subsection{Architectural Evolution}
Classical pipelines utilizing HRV, EDA, and ACC features combined with
SVM, RF, or kNN remain competitive in retrospective settings
\parencite{sethFeasibilityCardiacbasedSeizure2023,masonHeartRateVariability2024}.
However, deep temporal models such as 1D-CNN, LSTM, and TCN often
improve phase discrimination. Recently, attention variants and
Transformers are being explored on these signals
\parencite{yangMultimodalAISystem2022,pordoyEnhancedNonEEGMultimodal2025,
abtahiIdentificationRelevantECG2025,kalousiosECGbasedEpilepticSeizure2024}.
A key inquiry of this review is determining how feature-based models
compare to deep models and whether attention mechanisms add distinct
value under ambulatory, alarm-based evaluation.

\subsection{Fusion and Personalization}
Fusion across ECG, EDA, ACC, and PPG at the feature, decision, or
representation level seeks complementary cues and is relevant to
maintaining performance under noise and non-stationarity
\parencite{yangMultimodalAISystem2022,pordoyEnhancedNonEEGMultimodal2025}.
Additionally, patient-specific fine-tuning and post-hoc calibration
shift the operating point toward lower FAR/h and higher sensitivity in
ambulatory use \parencite{andradePerformanceSeizurePrediction2024}.
Thresholding and periodic recalibration are commonly used to handle
drift.

\subsection{Robustness and Generalization}
Robustness requires explicit checks of domain shift across patient,
device, site, and activity. Cross-dataset evaluation (e.g., training
on TUH and testing on RPAH or EPILEPSIAE) exposes this shift and favors
alarm-based over sample-based metrics
\parencite{yangMultimodalAISystem2022,andradePerformanceSeizurePrediction2024}.
Artifact handling, augmentation through time-warping, jitter, or
scaling, and Out-of-Distribution (OOD) detection support
generalization \parencite{kalousiosECGbasedEpilepticSeizure2024,
sethFeasibilityCardiacbasedSeizure2023}.

\subsection{Deployment Constraints}
Finally, deployment constraints bring attention to latency and energy
consumption. Compact sequence models and event-driven processing
reduce burden, while clear alarm rationales support safe clinical use.

\section*{Provisional Table of Contents (\textasciitilde15 Pages Text)}
\begin{enumerate}
  \item Introduction (Problem \& Motivation) -- 1.0 p
  \item Aim and Research Questions -- 0.5 p
  \item Methodology (Search, selection, extraction, concept matrix,
    PRISMA overview) -- 1.0 p
  \item Technical Background (Seizure taxonomy; autonomic markers;
    sensing accuracy) -- 1.0 p
  \item Modeling Landscape (Features vs. DL; CNN/LSTM/TCN;
    attention/transformers; fusion; personalization; robustness; alarm
    logic; deployment) -- 5--6 p
  \item Results (Concept matrix; comparative ranges; modality/model
    comparison; temporal \& personalization effects) -- 3--4 p
  \item Discussion (Trade-offs: sensitivity vs. FAR/h; robustness;
    interpretability; deployment feasibility) -- 1.5--2 p
  \item Limitations and Future Work -- 0.5--1 p
  \item Conclusion -- 0.5 p
\end{enumerate}